\documentclass[a4paper]{article}

\usepackage{anyfontsize}
\usepackage{amsmath,amssymb,mathrsfs}
\usepackage{autobreak}
\allowdisplaybreaks
\usepackage{bm}
\usepackage{indentfirst}
\setlength{\parindent}{0em}
\usepackage{parskip}
\usepackage{geometry}
\geometry{top=3cm,bottom=3cm,left=2.75cm,right=2.75cm}

\begin{document}

\title{\textbf{Exercise 2}}
\author{}
\date{}

\maketitle

\subsection*{Problem 1: Observational cosmology}

Consider a flat universe with $\Omega_\mathrm{m}=0.27$, $\Omega_\Lambda=0.73$, and $H=70\,\mathrm{km/(s\cdot Mpc)}$.

\textit{Hint:} Make numerical approximations whenever needed.

\begin{itemize}
    \item[(a)] We observe a quasar at redshift $z=3$ to have a brightness of $10^{-31}\,\mathrm{W/(m^2\cdot Hz)}$ at a wavelength of $600\,\mathrm{nm}$. To which intrinsic luminosity of the quasar in units of $\mathrm{W/Hz}$ at which wavelength does this observation correspond?
    
    \item[(b)] The quasar is located in a galaxy, whose diameter is 5 arcsec as seen by us. How big is this galaxy in units of physical kpc?
    
    \item[(c)] A gas cloud at the same redshift is located 10 arcsec away from the quasar brightness (in units of $\mathrm{W/(m^2\cdot Hz)}$) and at what wavelength does an observer within the gas cloud observe the intrinsic luminosity of the quasar of part (a)? Assume that the cloud is gravitationally bound to the galaxy, in which the quasar resides.
    
    \item[(d)] If the quasar is observed to change in brightness by us today, when would you expect us to see the gas cloud in part (c) respond to this change (assuming it responds instantaneously to changes in the incident flux)?
    
    \item[(e)] We observe the quasar to be part of a group of galaxies, i.e. the corresponding galaxies constitute a gravitationally bound system at the redshift of the quasar. What is the velocity dispersion $\sigma_v$ in $\mathrm{km/s}$ (i.e. the standard deviation of galaxy velocities in the restframe of the group), if we observe the galaxies to exhibit a standard deviation in redshift of $\sigma_z=0.001$?
    
    \item[(f)] Such quasars are believed to have a comoving number density of about $10^{-6}\,\mathrm{Mpc^{-3}}$
    i.e. they have a typical separation of about 100 comoving Mpc. How far away (i) in angle on the sky and (ii) in redshift, would you expect the next quasar to be, i.e. what angle $\Delta\theta$ and what $\Delta z$ correspond to 100 comoving Mpc at $z=3$?
\end{itemize}

\end{document}